\section*{Lab 7 (Forest Variants) Implementation Tasks}
\textbf{Lab 7 focus:} random partitioning as an anomaly signal (Isolation Forest), randomized split thresholds (Extra-Trees), and full leaf-distribution prediction intervals (Quantile Regression Forests).

\medskip
\textbf{Implementation policy (strict):}
\begin{itemize}
  \item Implement the isolation tree, the random-threshold split rule, and the quantile aggregation with NumPy only.
  \item Do not use a library implementation (e.g.\ \texttt{IsolationForest}, \texttt{ExtraTreesClassifier}) for the main task.
  \item You may reuse your decision tree implementation from Lab 5 and your random forest implementation from Lab 6.
  \item Plotting libraries are allowed for visual inspection.
  \item External libraries are allowed only for data loading if explicitly needed.
  \item Work through the notebook in order.
\end{itemize}

\section*{Core Equations And Rules}
\begin{itemize}
  \item Isolation tree split: pick feature $j$ uniformly at random, then pick threshold $t$ uniformly at random between the min and max of feature $j$ at the current node.
  \item Path length $h(x)$: number of edges from the root to the leaf containing $x$.
  \item Average path length normalizer over $\psi$ nodes:
  \[
    c(\psi) = 2H(\psi-1) - \frac{2(\psi-1)}{\psi},
    \qquad H(i) \approx \ln(i) + 0.5772.
  \]
  \item Anomaly score:
  \[
    s(x,\psi) = 2^{-\,\mathbb{E}[h(x)]/c(\psi)}.
  \]
  \item Extra-Trees split rule: for each candidate feature, draw \emph{one} threshold uniformly at random (no exhaustive search); keep the feature with the best $\Delta I$ among those random draws.
  \item QRF leaf weight for training example $i$:
  \[
    w_i(x) = \frac{1}{B}\sum_{b=1}^{B}
    \frac{\mathbf{1}\{x_i \in L_b(x)\}}{\bigl|L_b(x)\bigr|}.
  \]
  \item Weighted empirical CDF and quantile prediction:
  \[
    \hat{F}(y\mid x) = \sum_{i=1}^{n} w_i(x)\,\mathbf{1}\{y_i \le y\},
    \qquad
    \hat{Q}_\tau(x) = \inf\{y : \hat{F}(y\mid x) \ge \tau\}.
  \]
\end{itemize}

\section*{Data Files And Preparation}
\begin{itemize}
  \item \texttt{forest\_variants.ipynb}: guided notebook with the assignment steps.
  \item \texttt{advanced\_ai/forest/isolation\_forest.py}: starter implementation (Isolation Forest).
  \item \texttt{advanced\_ai/forest/extra\_trees.py}: starter implementation (Extra-Trees split rule).
  \item \texttt{advanced\_ai/forest/quantile\_forest.py}: starter implementation (Quantile Regression Forest aggregation).
  \item \texttt{prepare\_datasets.py}: generates the prepared data bundles.
  \item \texttt{scripts/prepare\_all\_data.sh}: shell entrypoint for dataset preparation.
  \item \texttt{scripts/check\_prepare\_all\_data.sh}: verification script for the prepared bundles.
  \item \texttt{toy\_isolation.py}, \texttt{extra\_trees\_iris.py}, \texttt{quantile\_forest\_demo.py}: task runners.
\end{itemize}

\medskip
Expected prepared files:
\begin{itemize}
  \item \texttt{datasets/isolation\_toy.npz}
  \item \texttt{datasets/extra\_trees\_iris.npz}
  \item \texttt{datasets/quantile\_regression.npz}
\end{itemize}

\medskip
The scripts use the conda environment \texttt{ai\_courses}.

\section*{Roadmap}
\begin{itemize}
  \item Start with Step 1: prepare and verify the datasets.
  \item After that, continue through the steps in order.
  \item Each step introduces only the task you need next.
\end{itemize}

\section*{Step 1: Prepare the Datasets}
\begin{itemize}
  \item From \texttt{assignment7\_forest\_variants/}, run:
  \[
    \texttt{bash scripts/prepare\_all\_data.sh}
  \]
  \item Then run:
  \[
    \texttt{bash scripts/check\_prepare\_all\_data.sh}
  \]
  \item Confirm that all three expected \texttt{.npz} bundles exist.
\end{itemize}

\section*{Step 2: Warm-up by Hand}
\begin{itemize}
  \item For $\psi=16$: $H(15) \approx 3.3182$. Compute $c(16)$ by hand.
  \item Point $A$ has $\mathbb{E}[h(A)] = 2.0$ and point $B$ has $\mathbb{E}[h(B)] = 6.0$. Compute $s(A,16)$ and $s(B,16)$, and state which point is the anomaly.
  \item A new point falls into leaves containing the pooled $y$-values (sorted)
  \[
    \{12,\,14,\,15,\,15,\,16,\,17,\,17,\,18,\,19,\,22\}.
  \]
  Compute the median ($\tau=0.5$) prediction and the $80\%$ prediction interval ($\tau=0.1$, $\tau=0.9$).
\end{itemize}

\section*{Step 3: Implement the Forest Variants Core}
\begin{enumerate}[label=\textbf{T\arabic*.}]
  \item \textbf{Isolation tree growth}
\begin{itemize}
  \item Open \texttt{advanced\_ai/forest/isolation\_forest.py}.
  \item Fill in \texttt{\_build\_node}.
  \item At each node, pick a feature $j$ uniformly at random, then a threshold $t$ uniformly at random between the min and max of feature $j$ at that node.
  \item Stop when a node has one point left, or the height limit $\approx \log_2 \psi$ is reached.
\end{itemize}

  \item \textbf{Path length}
\begin{itemize}
  \item Fill in \texttt{path\_length}.
  \item Route $x$ from the root and count the number of edges to its leaf.
\end{itemize}

  \item \textbf{Anomaly score}
\begin{itemize}
  \item Fill in \texttt{\_avg\_path\_length} (computes $c(\psi)$) and \texttt{anomaly\_score}.
  \item Average $h(x)$ over the $B$ isolation trees, then apply $s(x,\psi) = 2^{-\mathbb{E}[h(x)]/c(\psi)}$.
\end{itemize}

  \item \textbf{Extra-Trees split rule}
\begin{itemize}
  \item Open \texttt{advanced\_ai/forest/extra\_trees.py}.
  \item Fill in \texttt{\_random\_threshold\_split}.
  \item For each candidate feature, draw one threshold uniformly at random instead of searching all thresholds.
  \item Keep the feature whose random threshold gives the largest $\Delta I$.
  \item Reuse your Lab 5 impurity and Lab 6 bagging loop; only the threshold selection changes.
\end{itemize}

  \item \textbf{Leaf distribution collection}
\begin{itemize}
  \item Open \texttt{advanced\_ai/forest/quantile\_forest.py}.
  \item Fill in \texttt{\_collect\_leaf\_values}.
  \item For a regression forest trained as in Lab 6, store the training $y$-values that reach each leaf, for every tree.
\end{itemize}

  \item \textbf{Weighted quantile prediction}
\begin{itemize}
  \item Fill in \texttt{predict\_quantile}.
  \item Compute $w_i(x)$ from the leaves $x$ falls into across all $B$ trees.
  \item Build $\hat{F}(y\mid x)$ and return $\hat{Q}_\tau(x) = \inf\{y : \hat{F}(y\mid x) \ge \tau\}$.
\end{itemize}
\end{enumerate}

\section*{Step 4: Run the Self-Checks}
\begin{itemize}
  \item After implementing the core classes, run:
  \[
    \texttt{python -m advanced\_ai.self\_checks}
  \]
  \item Continue only after every check passes.
\end{itemize}

\section*{Step 5: Toy Isolation Forest}
\begin{itemize}
  \item Run \texttt{toy\_isolation.py} on the toy dataset with injected outliers.
  \item Report the average path length and anomaly score for an injected outlier and for a normal point.
  \item Save the anomaly score plot (score versus normalized path length, as in the lecture).
  \item State the threshold you used to flag anomalies and how many points were flagged.
\end{itemize}

\section*{Step 6: Extra-Trees vs Random Forest on Iris}
\begin{itemize}
  \item Run \texttt{extra\_trees\_iris.py}.
  \item Sweep over $B = 1, 5, 10, 20, 50, 100$ for both your Lab 6 random forest and your Extra-Trees forest.
  \item Record training time, training accuracy, and test accuracy for each $B$ and each method.
  \item State which method trains faster and which has lower variance across repeated runs.
\end{itemize}

\section*{Step 7: Quantile Regression Forest}
\begin{itemize}
  \item Run \texttt{quantile\_forest\_demo.py} on the 1D regression dataset.
  \item Report the median prediction and an $80\%$ prediction interval ($\tau=0.1,0.9$) for three test points.
  \item Save the plot of the predicted median curve with the prediction interval band.
  \item Report the fraction of held-out test points that actually fall inside their predicted $80\%$ interval.
\end{itemize}

\section*{Step 8: Optional Extension}
\begin{itemize}
  \item Vary the subsample size $\psi \in \{64, 128, 256, 512\}$ and report how the anomaly scores change.
  \item Try a $50\%$ and a $95\%$ prediction interval instead of $80\%$ and compare interval width.
  \item Combine Extra-Trees splitting with the quantile-forest leaf aggregation and report whether the prediction intervals change.
\end{itemize}

\section*{Submission}
\begin{itemize}
  \item Submit the notebook, the completed starter files, and the generated output files from the required runners.
  \item Use \texttt{collect\_submission.ipynb} or \texttt{collectSubmission.sh} to bundle the submission.
\end{itemize}
